\vspace{5mm}
\begin{ledgroup}
\footnotesize
\pstart
\noindent%
Die folgenden Stücke N.\;\ref{60308} und N.\;\ref{60307} sind inhaltlich und formal eng miteinander verbunden. 
Beim Konzept N.\;\ref{60307} handelt es sich um eine Ausarbeitung der kürzeren, hauptsächlich aus Rechnungen bestehenden Aufzeichnung N.\;\ref{60308}.
Die Berechnungen der vektoriellen Geschwindigkeiten und der Impulse beim Zweikörperstoß aus N.\;\ref{60308} 
kommen, trotz abweichender Nummerierung, in N.\;\ref{60307} unter Verwendung ähnlicher Kennzeichen wieder vor. 
Dabei weitet Leibniz einerseits die Betrachtung auf die Berechnung einer dritten Größe, der quadratischen \textit{vis}, aus;
andererseits ergänzt er die Rechnungen um einige methodologische Bemerkungen (z.\,B. über die dimensional aufsteigende Hierarchie der genannten physikalischen Größen und deren Erhaltungssätze)
und um eine kurze einleitende Erläuterung der verwendeten Symbolik.
An letztgenannter Stelle von N.\;\ref{60307} wird das Abhängigkeitsverhältnis beider Stücke deutlich:
Während das Konzept zunächst von einer Gleichung seinen Ausgang nahm, und zwar derselben wie bereits die Aufzeichnung N.\;\ref{60308},
hat Leibniz diesen Anfang gestrichen und den Gleichungen bzw. \glqq Theoremen\grqq\ die genannte Einleitung vorhergehen lassen (siehe die Varianten zu \refpassage{37_05_062_1a}{37_05_062_1b}).
\pend
\pstart
Das \href{\Scan{LH_37_05_062wz}}%
{Wasserzeichen} im Träger von N.\;\ref{60307} (einfaches Posthorn ohne Beizeichen) 
konnte noch nicht identifiziert werden.
Vielmehr stützt sich die vorgeschlagene Datierung auf Bezüge zu weiteren Texten: 
einerseits auf die terminologischen Ähnlichkeiten zwischen den vorliegenden Stücken und einem auf den Zeitraum 1696\textendash1698 datierbaren Text (\cite{02147}\textit{LSB} VIII,~3 N.\;76),
andererseits auf den Umstand, dass die besonders in N.\;\ref{60307} ausgearbeitete Liste dreier Erhaltungssätze samt ihrer dimensional aufsteigenden Anordnung (\refpassage{37_05_062_7a}{37_05_062_7b}) 
ein Kernstück des etwa 1700 entstandenen, unveröffentlichten \cite{02181}\textit{Essay de Dynamique} bildet (\textit{LMG} VI, S.\;215\textendash231, bes. S.\;226\textendash228).
\pend
\pstart
Bei dem ersten Text handelt es sich um 
\cite{02147}kommentierte Auszüge aus
\indexN{\textls{Sohier}, G.~J. gest. nach 1696}G.\,J.\,\textsc{Sohier}, 
\cite{00526}\glqq Nouveau systeme de la percussion des corps\grqq,
\cite{00157}\textit{JS} 7.~Mai 1696, S.\;211\textendash214.
Erstens verwendet Leibniz in den Auszügen dieselbe (ins Französische übertragene) Terminologie der \textit{vires impeditivae} bzw. \textit{progressiones} wie in N.\;\ref{60308} und N.\;\ref{60307} (\textit{LSB} VIII,~3, S.\;792\,f.). 
Zweitens schreibt er in einer gestrichenen Variante einige lateinische Sätze, die mit dem letzten Absatz von N.\;\ref{60307} (\glqq Notatu dignum \lbrack...\rbrack\ cum aequalis progressio aggressioni\grqq, hier auf \refpassage{37_05_062_6a}{37_05_062_6b}) nahezu wörtlich übereinstimmen: 
"Notandum autem quantum corpus corpori dat progressionis suae, tantum amittit. 
Unde paradoxum, omnibus redactis ad quietem servari posse quantitatem progressionis totalem, 
ea enim nulla erit si quantum erat progressus tantum obgressus"
(Variante \textit{(1)} zu S.\;793 Z.~18\textendash S.\;794 Z.~1).
Diese auffallende Übereinstimmung ist vielleicht dadurch zu erklären, dass Leibniz beim Verfassen der Auszüge die ihm vorliegenden Sätze aus N.\;\ref{60308} in den französischen Text eingearbeitet hat.
\pend
\pstart
Die Begrifflichkeit der \glqq force impeditive ou progressive\grqq\ kommt außerdem 
bei der Behandlung von Stoßfragen im Briefwechsel mit Papin der Jahre 1696\textendash1698 häufig vor: 
siehe dazu die editorische Datierung der \indexN{\textls{Sohier}, G.~J. gest. nach 1696}Sohier-Auszüge (\textit{LSB} VIII,~3 N.\;76, S.\;790\,f.) und Leibnizens Briefe an Papin
\cite{02071}vom 16. (26.) Juli 1696, 
\cite{01503}8. (18.) November 1697 und 
\cite{02073}16. (26.) Januar 1698
(\textit{LSB} III, 7 N.\;8, N.\;156 resp. N.\;177).
Aus den genannten Gründen wird für N.\;\ref{60308} und N.\;\ref{60307}
die Datierungsspanne 1696 bis 1700 vorgeschlagen.
\pend 
\end{ledgroup}
